\documentclass{ctexart}

\usepackage{xcolor}
\usepackage{fontspec}
\usepackage{ruby} 

\usepackage{xcolor}
\usepackage{xparse}
\usepackage{calc}

\usepackage{relsize}

\usepackage{etoolbox,xstring}

\usepackage{geometry}
\geometry{legalpaper, portrait, margin=2cm}

\definecolor{primary}{RGB}{47, 128, 190}    % Muted Blue
\definecolor{secondary}{RGB}{192, 57, 43}   % Deep Red
\definecolor{accent}{RGB}{39, 174, 96}      % Cool Green
\definecolor{background}{RGB}{255, 255, 255} % Softer Light Gray
\definecolor{link}{RGB}{128, 0, 128}       % Regal Purple

\definecolor{pagedark}{RGB}{0, 0, 0}            %black
\definecolor{textdark}{RGB}{234, 234, 234}      %grey

\definecolor{pagelight}{RGB}{255, 255, 255}            %white
\definecolor{textlight}{RGB}{0, 0, 0}            %black

\colorlet{page}{pagedark}
\colorlet{text}{textdark}

\pagecolor{page}
\color{text}

\setmainfont{NotoSans}[
  Path               = fonts/,
  Extension          = .ttf,
  UprightFont        = *-Light,
  ItalicFont         = *-LightItalic,
  BoldFont           = *-Bold,
  BoldItalicFont     = *-BoldItalic,
]
\newcommand\textrmlf[1]{{\NHLight#1}}

\setCJKmainfont{NotoSansJP}[
  Path               = fonts/,
  Extension          = .ttf,
  UprightFont        = *-Light,
  BoldFont           = *-Bold,
]

\newfontfamily\textserif{NotoSerif}[
  Path               = fonts/,
  Extension          = .ttf,
  UprightFont        = *-Light,
  ItalicFont         = *-LightItalic,
  BoldFont           = *-Bold,
  BoldItalicFont     = *-BoldItalic,
]

\newCJKfontfamily\textCJKserif{NotoSerifJP}[
  Path               = fonts/,
  Extension          = .ttf,
  UprightFont        = *-Light,
  BoldFont           = *-Bold,
]

\newfontfamily\textfancy{NotoSerif}[
  Path               = fonts/,
  Extension          = .ttf,
  UprightFont        = *-Light,
  ItalicFont         = *-LightItalic,
  BoldFont           = *-Bold,
  BoldItalicFont     = *-BoldItalic,
]

\newfontfamily\textcharm{EagleLake}[
  Path               = fonts/,
  Extension          = .ttf,
  UprightFont        = *-Regular,
]

\newCJKfontfamily\textCJKfancy{KaiseiHarunoUmi}[
  Path               = fonts/,
  Extension          = .ttf,
  UprightFont        = *-Regular,
  BoldFont           = *-Bold,
]

% define \cjruby{漢字}{かんじ}
\NewDocumentCommand{\cjruby}{mm}{%
  {%
    \ruby{\textCJKserif{}#1}{#2}
  }%
}


%\NewDocumentCommand{\frby}[2][1]{
%    \ruby{\textCJKfancy{}#1}{#2}
%}

\NewDocumentCommand{\internalfrby}{m m}{%
  \ruby{\textCJKfancy{}#1}{\smaller\smaller\textcharm{}#2}
}

\NewDocumentCommand{\frby}{m O{}}{%
  \vphantom{\makebox[0pt][l]{\internalfrby{語}{M}}}%
  \internalfrby{#1}{#2}
}

\ExplSyntaxOn
\NewDocumentCommand\fruby{>{\SplitList{~}}m}{\ProcessList{#1}{\fruby_map}}
\NewDocumentCommand\fruby_map{>{\SplitArgument{1}{=}}m}{{\fruby_impl #1}}
\NewDocumentCommand\fruby_impl{m m}{\IfNoValueTF{#2}{\frby{#1}}{\frby{#1}[#2]}}
\ExplSyntaxOff

\def\maxvalue(#1,#2){\ifnum #1>#2 #1\else #2\fi}

\NewDocumentCommand{\englishundertokisin}{m m}{%
  \begin{minipage}[c]{\maxof{\widthof{#1}}{\widthof{#2}}}
    \centering{#1}\small\\\centering{#2}
  \end{minipage}%
}


\begin{document}

\begin{center}
    \vspace*{2cm}
    {\fontsize{48}{60}\bfseries\color{primary}\fruby{獣=akisi あの 狸=sowili}}
    \vspace{1cm}
\end{center}
\vspace{1cm}
\hrule
\vspace{1cm}

\section[人去ら]{\Huge\fruby{人 去 ら}}
{
\noindent
\begin{minipage}[t]{3.25cm}
{\textCJKfancy{}\LARGE
  \fruby{人 去 ら ：}\\
  \fruby{り 去 市=ison ・}\\
  \fruby{り 市=ison い 食=moko ・}\\
  \fruby{り 開 去 ／}\\
}
\end{minipage}
\begin{minipage}[t]{3.25cm}
{\textCJKfancy{}\LARGE
  \fruby{り 終=pini 多=moti ・}\\
  \fruby{り 丸=siki 切=kipisi 、}\\
  \fruby{于 音=kalama 大=soli ・}\\
  \fruby{于 獣=akisi 大=soli ！}\\
}
\end{minipage}
\begin{minipage}[t]{3.75cm}
{\textCJKfancy{}\LARGE
  \fruby{獣=akisi ら ：}\\
  \fruby{り 黄=silo 木=kasi 、}\\
  \fruby{り 大=soli 大=soli ！}\\
  \fruby{り 有=jo い 空=kon 悪=iki ！}\\
}
\end{minipage}
}

\section[獣ろ鬥い人]{\Huge\fruby{獣=akisi ろ 鬥=otala い 人}}

\noindent
\begin{minipage}[t]{5.25cm}
{\textCJKfancy{}\LARGE
  \fruby{獣=akisi ろ 黄=silo 木=kasi ら 、}
  \fruby{り 要=wili ／ り 要=wili 食=moko 、}\\
  \fruby{い 食=moko ろ い 市=ison 、}\\
  \fruby{い 人=jan ろ 市=ison 、}\\
  \fruby{り 要=wili 食=moko い 家=tomo 市=ison ！}\\
}
\end{minipage}
\begin{minipage}[t]{3.75cm}
{\textCJKfancy{}\LARGE
  \fruby{人 去=tawa り 去=tawa 許=taso 不=ala}\\
  \fruby{人 去=tawa り 人 鬥=otala ！}\\
  \fruby{人 鬥=otala ろ 去=tawa ら}\\
  \fruby{り 開=opin 有=jo い 具=ilo 鬥=otala}\\
  \fruby{り 作=pali い 穴=lopa ひ 内 獣=akisi 、}\\
  \fruby{獣=akisi ろ 黄=silo 木=kasi 、}\\
}
\end{minipage}

\newpage
\setcounter{section}{0}
\begin{center}
    \vspace*{2cm}
    {\fontsize{48}{60}\bfseries\color{primary}\fruby{獣=akisi あの 狸=sowili}}\LARGE\\
    {\Huge\textserif{}Akisi or Sowili?}
    \vspace{1cm}
\end{center}
\vspace{1cm}
\hrule
\vspace{1cm}

\section[人去ら]{
\Huge\englishundertokisin
  {\fruby{人 去 ら}}
  {\large\normalfont{}The Traveler}
}
{
\noindent
\begin{minipage}[t]{3.25cm}
{\textCJKfancy{}\LARGE
  \fruby{人 去 ら ：}\\
  \fruby{り 去 市=ison ・}\\
  \fruby{り 市=ison い 食=moko ・}\\
  \fruby{り 開 去 ／}\\
  \scriptsize
  There was a traveler;\\
  she entered the market,\\
  she bought some food,\\
  she began to leave...\\
}
\end{minipage}
\begin{minipage}[t]{3.25cm}
{\textCJKfancy{}\LARGE
  \fruby{り 終=pini 多=moti ・}\\
  \fruby{り 丸=siki 切=kipisi 、}\\
  \fruby{于 音=kalama 大=soli ・}\\
  \fruby{于 獣=akisi 大=soli ！}\\
  \scriptsize
  she stopped suddenly.\\
  She turned around,\\
  There was a loud sound.\\
  There was a huge creature!\\
}
\end{minipage}
\begin{minipage}[t]{3.75cm}
{\textCJKfancy{}\LARGE
  \fruby{獣=akisi ら ：}\\
  \fruby{り 黄=jelo 木=kasi 、}\\
  \fruby{り 大=soli 大=soli ！}\\
  \fruby{り 有=jo い 空=kon 悪=iki ！}\\
  \scriptsize
  This creature, \\
  it was green like a tree,\\
  it was very large!\\
  It had an evil spirit!\\
}
\end{minipage}
}

\section[獣ろ黄木]{
  \Huge\englishundertokisin
    {\fruby{獣=akisi ろ 黄=jelo 木=kasi}}
    {\large\normalfont{}The Green Creature}
}

\noindent
\begin{minipage}[t]{5.25cm}
{\textCJKfancy{}\LARGE
  \fruby{獣=akisi ろ 黄=jelo 木=kasi ら 、}
  \fruby{り 要=wili ／ り 要=wili 食=moko 、}\\
  \fruby{い 食=moko ろ い 市=ison 、}\\
  \fruby{い 人=jan ろ 市=ison 、}\\
  \fruby{り 要=wili 食=moko い 家=tomo 市=ison ！}\\
  \scriptsize
  This creature, green like a tree,\\
  it wanted... it wanted to eat,\\
  to eat the market food,\\
  to eat the market people.\\
  It wanted to eat the market house!\\
}
\end{minipage}

\newpage
% Introduction Section
\section[Appendix]{\Huge{}Appendix}
This is written in 言新 {\it{}(toki sin)}.\\

Here are the 画多 {\it{}(sitilin moti)},
common 画字 {\it{}(sitilin kansi)}.
画 {\it{}(sitilin)} which aren't on this list will frequently have ruby text above them to assist in reading.\\

\Huge
\fruby{私=mi 君=sina 其=ona}

\fruby{因=tan 去=tawa 要=wili 集=alasa 待=awin 来=kama 能=kin}

\fruby{ら=la り=li ろ=lo}

\fruby{い=i か=ka}

\fruby{いん=in あの=ano ひ=pi 許=taso 内=insa}

\fruby{不=ala}

\fruby{人=jan 于=lon 去=tawa}



\newpage



\end{document}
